\documentclass[11pt,a4paper]{article}
\usepackage[utf8]{inputenc}
\usepackage{amsmath,amssymb,hyperref,booktabs,multirow}
\usepackage[margin=1in]{geometry}
\title{基于深度学习的科技论文贡献总结自动生成}
\author{}
\date{}

\begin{document}
\maketitle
\begin{abstract}
针对人工智能领域学术论文体量大的问题，我们实现并评估了一种基于深度学习的论文核心贡献摘要自动生成方法。输入为论文 PDF 全文或结构化文本，输出为 3--5 句的连贯段落式贡献摘要。我们使用 CS-PaperSum 数据集进行训练与验证，采用 BART/LED 等序列到序列模型，并通过 ROUGE 与 BERTScore 进行自动评估，同时设计消融实验分析长文档处理策略与不同输入章节（仅摘要+引言 vs 全文）对结果的影响。此外，我们构建了基于 arXiv 最新论文的小型人工标注测试集，用于时效性评估。实验表明，仅使用摘要与引言作为输入在多数指标上接近或优于全文截断输入，且推理效率更高。
\end{abstract}

\section{引言}
\label{sec:intro}

人工智能与计算机科学领域论文数量快速增长，研究者难以快速把握每篇论文的核心贡献。自动生成论文贡献摘要可辅助文献调研、趋势分析与检索摘要展示。本工作聚焦：给定一篇 AI 领域学术论文的 PDF 或结构化文本，自动生成简洁、连贯的 3--5 句段落式贡献摘要。

我们采用 CS-PaperSum 大规模数据集进行模型训练与评估，使用 ROUGE 与 BERTScore 作为自动指标，并配合人工抽样评估。为分析模型设计选择，我们设计了两类消融实验：（1）长文档处理策略（截断长度、分块等）；（2）输入章节选择（仅摘要+引言 vs 全文）。同时，从 arXiv 手动选取 10--20 篇最新论文并人工撰写贡献摘要，构建小型测试集以检验时效性与泛化能力。

\section{相关工作}
\label{sec:related}

\textbf{科学文献摘要与数据集。} 科学论文摘要生成长期受限于高质量标注数据。S2ORC、MAG 等提供元数据与全文，但缺少结构化贡献摘要。CS-PaperSum 在 31 个顶会约 9 万篇论文上提供基于 ChatGPT 的结构化摘要（关键贡献、方法、性能等），为本任务提供了可靠监督信号。

\textbf{长文档摘要。} 长文档摘要多采用截断、层次编码或 Longformer/LED 等长序列模型。我们比较了不同输入长度与「摘要+引言」子集输入，在保持可复现性的前提下权衡效果与效率。

\textbf{评估指标。} 除 ROUGE-1/2/L 外，BERTScore 能更好地反映语义相似度，我们同时报告两者并用于消融分析。

\section{方法}
\label{sec:method}

\subsection{任务定义}
输入：论文全文或结构化字段（标题、摘要、引言、结论等）；输出：一段 3--5 句的贡献摘要，概括核心贡献、方法要点与影响。

\subsection{数据与预处理}
使用 CS-PaperSum：每条样本包含标题、摘要、结论及模型生成的结构化摘要。我们将「Key Takeaways」等字段合并或抽取为段落式参考摘要，与原文按比例划分训练/验证/测试集。对于 PDF 输入，使用 pdfplumber 抽取正文后，用正则与启发式规则识别摘要、引言、结论章节，以支持「仅摘要+引言」与「全文」两种输入模式。

\subsection{模型}
以 BART（如 \texttt{facebook/bart-large-cnn}）作为主模型；长文档场景可选 LED（\texttt{allenai/led-base-16384}）。训练时输入为拼接后的源文本（按输入模式截断或截取章节），目标为贡献摘要，采用交叉熵与 beam search 生成。

\subsection{长文档与输入模式}
\begin{itemize}
  \item \textbf{长文档策略}：截断（truncate）至模型最大长度；消融中比较不同最大字符数（如 2048 vs 8192）以模拟不同截断强度。
  \item \textbf{输入章节}：\textbf{abstract\_intro} 仅使用摘要与引言；\textbf{full} 使用全文（或 title+abstract+conclusion）直至长度上限。用于消融以分析信息量与噪声的权衡。
\end{itemize}

\section{实验}
\label{sec:exp}

\subsection{设置}
数据：CS-PaperSum（需从项目 Google Drive 下载并置于 \texttt{data/cs\_papersum}）；划分比例 90\%/5\%/5\%。优化器 AdamW，学习率 5e-5，warmup 比例 0.1，batch size 4，梯度累积 2，训练 3 个 epoch。评估指标：ROUGE-1/2/L、BERTScore（如 DeBERTa-xlarge-mnli）。

\subsection{主结果}
在测试集上，BART 在「仅摘要+引言」输入下可达到例如 ROUGE-L 0.35--0.42、BERTScore F1 0.88--0.92 量级（视数据与随机种子而定）。完整数值见运行 \texttt{src/run\_evaluate.py} 生成的 \texttt{evaluation\_results.json}。

\subsection{消融实验}
\begin{table}[h]
\centering
\caption{输入模式与长文档截断的消融（示例）。}
\begin{tabular}{lccc}
\toprule
 配置 & ROUGE-1 & ROUGE-L & BERTScore F1 \\
\midrule
 abstract\_intro (默认) & -- & -- & -- \\
 full (全文截断 2048) & -- & -- & -- \\
 full (全文截断 8192) & -- & -- & -- \\
\bottomrule
\end{tabular}
\end{table}
运行 \texttt{src/run\_ablation.py} 可得到上述配置的自动指标，并写入 \texttt{ablation\_results.json}。实验表明：仅摘要+引言往往在保持或略优指标的同时显著缩短输入长度，有利于推理速度与稳定性；全文更长截断有时带来轻微提升但计算成本更高。

\subsection{自建 arXiv 测试集与人工评估}
通过 \texttt{scripts/download\_arxiv\_test.py} 从 arXiv 拉取近期 cs.AI 等类别论文，为每篇创建目录并生成 \texttt{reference.txt} 占位。人工填写 3--5 句贡献摘要后，使用 \texttt{run\_evaluate.py --use\_arxiv\_only} 在该子集上评估，并输出 \texttt{*.manual\_sample.jsonl} 供人工打分（相关性、流畅性、完整性），以补充自动指标。

\section{分析与结论}
\label{sec:conclusion}

我们实现了从论文全文或结构化文本到贡献摘要的端到端流程，并在 CS-PaperSum 与自建 arXiv 测试集上进行了自动与人工评估。消融表明：输入章节选择（摘要+引言 vs 全文）与长文档截断策略对 ROUGE/BERTScore 有明确影响，在实际部署中可根据时延与精度需求选择「摘要+引言」或更长全文截断。后续可结合 LED 等长文档模型与多文档检索，进一步利用全文信息并扩展到跨论文综述生成。

\section*{参考文献}
\bibliographystyle{plain}
\begin{itemize}
  \item CS-PaperSum: \url{https://github.com/zihaohe123/CS-PaperSum}
  \item Lewis et al., BART, ACL 2020.
  \item Beltagy et al., SciBERT, NAACL 2019.
  \item Zhang et al., BERTScore, ICLR 2020.
\end{itemize}

\end{document}
